% SPDX-License-Identifier: CC-BY-SA-4.0
% Author: Matthieu Perrin

\begingroup


\SetKwFunction{Add}{add}
\SetKwFunction{Insert}{insert}
\SetKwFunction{Delete}{delete}
\SetKwFunction{Contains}{contains}
\SetKwData{Inserted}{inserted}
\SetKwData{Deleted}{deleted}

\begin{exercice}[L'ensemble OR-Set]

  Le but de cet exercice est de concevoir un ensemble permettant d'insérer et de supprimer des éléments,
  tout en garantissant l'\emph{eventual consistency}.

  On considère d'abord une machine à états dont l'état est constitué de deux listes :
  \begin{itemize}
    \item la liste \Inserted, contenant toutes les valeurs insérées, et
    \item la liste \Deleted, contenant toutes les valeurs supprimées.
  \end{itemize}

  Les opérations sont :
  \begin{itemize}
  \item $\Insert(v): \Inserted.\Add(v)$, 
  \item $\Delete(v): \Deleted.\Add(v)$,
  \item $\Contains(v): \Return (v \in \Inserted \setminus \Deleted)$.
  \end{itemize}

  \begin{question}
    \item Montrer que toutes les opérations de mise à jour commutent. Que peut-on en déduire ?
  \end{question}

  Le problème de la machine à états précédente est qu'il est impossible de réinsérer une valeur après sa suppression.  
  Pour résoudre ce problème, on définit l'\emph{OR-Set} (\emph{Observed-Remove Set}).

  \begin{itemize}
  \item Lorsqu'on exécute $\Insert(v)$, on génère un identifiant $t$ unique (par exemple une estampille de Lamport) et on insère le couple $\langle v, t\rangle$ dans la liste \Inserted.  
  \item Lorsqu'on exécute $\Delete(v)$, on ajoute à \Deleted\ tous les couples $\langle v, t\rangle$ actuellement présents dans \Inserted.  
  \item L'opération $\Contains(v)$ renvoie vrai s'il reste au moins un couple $\langle v, t\rangle$ dans $\Inserted \setminus \Deleted$.
  \end{itemize}
  
  \begin{question}
    \item Écrire l'algorithme de l'OR-Set.
    \item Modifier l'algorithme pour réduire au maximum la complexité spatiale.
      On pourra utiliser la diffusion causale.
  \end{question}

\end{exercice}

\endgroup
\endinput
